\documentclass[12pt]{article}
\usepackage[utf8]{inputenc}
\usepackage[T1]{fontenc}
\usepackage{hyperref}
\title{Malagasy, Tandroy-Mahafaly Language Resources for K-12 Curriculum Development}
\author{Compiled by Claude Sonnet 5 (claude-sonnet-5) for Andr\'as Kornai}
\date{}
\begin{document}
\maketitle

\section{Introduction}
Tandroy-Mahafaly Malagasy (ISO 639-3 tdx; Glottocode tand1256) is a variety of the Malagasy macrolanguage spoken in southern Madagascar. Karimbola is considered a dialect of Tandroy, and Mahafaly is a closely related variety. Specific Mahafaly resources were not found.

\section{Primary Educational Resources}
No Tandroy primers or textbooks were found in the searches.

\section{Linguistic and Reference Materials}
\begin{itemize}
\item Wikipedia, Tandroy dialect: \url{https://en.wikipedia.org/wiki/Tandroy_dialect}
\item Omniglot page on Tandroy: \url{https://www.omniglot.com/babel/tandroy.htm}
\item Bilingual Malagasy dialect--French dictionary (Tandroy section, 1996, with an introduction to Malagasy grammar; authors include Narivelo Rajaonarimanana): \url{https://asiatheque.com/en/book/dictionnaire-malgache-dialectal-francais}
\end{itemize}

\section{Religious and Community Resources}
\begin{itemize}
\item Tandroy Bible, \textit{Saontsin' A\~nahare}, digital edition (marked as a draft, Creative Commons): \url{https://ebible.org/details.php?id=tdx}
\item Report on the translation of the complete Bible into Tandroy: \url{https://www.cgmmag.com/?p=4690}
\item Word und Tat, Madagascar page: \url{https://www.wortundtat.de/en/madagascar/}
\item Joshua Project, Antandroy: \url{https://joshuaproject.net/people_groups/10345/MA}; language page: \url{https://www.joshuaproject.net/index.php/languages/tdx}. Its Scripture data appears out of date.
\end{itemize}

\section{Digital and Audio Resources}
The eBible edition above is the only digital text found. No audio or corpus was seen.

\section{Curriculum Development Context}
The complete Bible gives a written Tandroy text base and an orthography in use. Reference material for the language is thin, and no school materials were found. The Wikipedia, Omniglot and Word und Tat pages were not read in detail.
\end{document}